\documentclass[pdflatex,sn-mathphys-num]{sn-jnl}

\usepackage{graphicx}
\usepackage{multirow}
\usepackage{amsmath,amssymb,amsfonts}
\usepackage{amsthm}
% \usepackage{mathrsfs}
\usepackage{lmodern}
\usepackage[title]{appendix}
\usepackage{xcolor}
\usepackage{textcomp}
\usepackage{manyfoot}
\usepackage{booktabs}
\usepackage{algorithm}
\usepackage{algorithmicx}
\usepackage{algpseudocode}
\usepackage{listings}
\usepackage{array}
\usepackage{makecell}
\usepackage{url}

\raggedbottom

\begin{document}

\title[Monocular Non-rigid Endoscopic 3D Reconstruction]{Monocular Non-rigid Endoscopic 3D Reconstruction with Tissue-aware Heterogeneous Physical Priors}

% Replace the author and affiliation information below with the final submission details.
\author*[1]{\fnm{Author} \sur{Name}}\email{author@example.com}
\affil*[1]{\orgdiv{Department}, \orgname{Institution}, \orgaddress{\city{City}, \country{Country}}}

\abstract{\textbf{Purpose:} Minimally invasive surgery benefits from small incisions and reduced patient trauma, but reliable intraoperative geometric perception remains essential for image-guided surgical applications. Monocular non-rigid SLAM can recover camera motion and sparse intraoperative structure, yet sparse surfaces usually provide limited geometric coverage, and dense depth completion may introduce affine scale-and-shift mismatch and cross-frame geometric inconsistency. To address these reconstruction-level limitations, we propose a monocular non-rigid endoscopic reconstruction method with tissue-aware heterogeneous physical priors.

\textbf{Methods:} In the front end, a tissue-aware heterogeneous physical graph (THPG) is constructed from tissue-region priors and node-level relative physical weights to support tracking and sparse geometric recovery. Unlike methods that use uniform global viscoelastic coefficients, the proposed front end assigns heterogeneous elastic and viscous weights to graph nodes and edges. In the back end, THPG topology, sparse physical anchors, and tracking states estimated by the front end are propagated as structured guidance for test-time diffusion-based depth completion, with the goal of densifying the intraoperative surface while improving pseudo-metric alignment to sparse physical anchors and local geometric consistency across frames. After multi-frame fusion, an auxiliary THPG-guided anchor-preserving refinement step is used to suppress residual high-frequency noise while limiting changes to sparse-anchor alignment and deformation-aware local topology.

\textbf{Results:} Experiments on public endoscopic datasets and a self-collected deformable organ phantom dataset show favorable tracking and dense depth accuracy under the tested settings. In the front-end evaluation, the proposed method obtained RMSE values of 1.66~mm, 0.76~mm, and 0.42~mm on Hamlyn abdominal-sequence-6, exploration-sequence-20, and liver-seq-21, respectively. In the back-end dense depth evaluation, the RMSE values were 2.31 on Hamlyn exploration-sequence-20 and 2.76 on SCARED Dataset~1, Sequence~2. Together with the ablation results reported in the main text, these observations suggest improved compatibility between dense recovery and front-end pseudo-metric anchors within the evaluated sequences.

\textbf{Conclusion:} The observed results indicate that explicit cross-stage propagation of tissue-aware heterogeneous physical priors is associated with lower front-end tracking and back-end dense reconstruction errors under the evaluated conditions, without retraining the foundation depth model. The method provides a denser and more scale-consistent intraoperative point cloud representation, while direct validation in downstream surgical-navigation tasks remains a topic for future work.}

\keywords{monocular endoscopy, non-rigid 3D reconstruction, scale-consistent dense reconstruction, tissue-aware heterogeneous physical graph, diffusion depth completion, multi-frame fusion}

\maketitle

\section{Introduction}\label{sec:introduction}

Minimally invasive surgery has the advantages of small incisions, reduced tissue trauma, faster postoperative recovery, and lower patient burden. To fully exploit these advantages in image-guided surgery, the intraoperative scene must be reconstructed in a geometrically reliable form. Dense and scale-consistent intraoperative geometry can provide useful visual and geometric support for subsequent surgical navigation tasks, including surface inspection, surgical field measurement, augmented reality guidance, and potential point-cloud-based alignment with preoperative models \cite{bib1,bib2}. Therefore, monocular endoscopic 3D reconstruction is not only a problem of visualizing the surgical field, but also a prerequisite for obtaining reliable intraoperative geometric representations.

In conventional image-guided endoscopic minimally invasive surgery, endoscope localization often relies on external optical tracking systems. Such solutions require additional hardware and calibration procedures, which increase system cost and deployment complexity. Moreover, they typically depend on maintaining a fixed spatial relationship between the endoscope camera and the tracked marker assembly, making them more naturally compatible with rigid endoscopes and less convenient for flexible or more freely manipulated imaging setups. In contrast, monocular vision-based reconstruction offers the possibility of recovering camera motion and scene structure directly from the endoscopic video itself, thereby reducing reliance on dedicated optical tracking hardware and relaxing the need for rigid tracker--camera coupling.

Despite this potential, endoscopic scenes remain particularly challenging. Tissue surfaces undergo persistent non-rigid deformation caused by respiration, heartbeat, and instrument interaction. Such continuous soft-tissue motion directly violates the static-scene assumption underlying conventional SLAM and multi-view geometry, often leading to pose drift and geometric distortion \cite{bib1,bib3,bib4}. Moreover, many SLAM-based methods naturally produce sparse map points or sparse surface structures. Although such sparse maps can support camera tracking, they usually provide limited surface coverage and sparse local neighborhoods, which restricts their usefulness for dense geometric analysis and for potential downstream correspondence-based applications.

Dense recovery is therefore needed to convert the sparse intraoperative map into a more complete surface representation. However, densification alone is not sufficient. In monocular endoscopy, both front-end tracking and back-end dense depth prediction lack direct absolute scale observations. If the dense point cloud generated for each frame is not aligned to the same pseudo-metric coordinate convention as the front-end sparse anchors, the same tissue surface may appear with different depth scales or affine shifts across frames. Such scale mismatch can be accumulated during multi-frame fusion and may be incorrectly interpreted as tissue deformation or structural inconsistency in subsequent geometric analysis. Consequently, the key requirement is not only to generate a denser point cloud, but also to keep the dense reconstruction geometrically compatible with the front-end sparse structure in affine scale and shift as well as local surface topology.

Existing studies address these issues only partially. Deformation-aware SLAM methods, such as DefSLAM, SD-DefSLAM, DSDT, MCPD, and NR-SLAM, improve pose estimation and sparse reconstruction by introducing deformation modeling, local co-motion constraints, and dynamic map maintenance \cite{bib5,bib6,bib7,bib8,bib9}. However, these methods mainly use the estimated deformation model to stabilize the front end, and the resulting structure remains sparse for downstream dense geometric use. Meanwhile, depth completion and generative monocular depth recovery offer an alternative route toward dense reconstruction \cite{bib10,bib11,bib12,bib13,bib14,bib15,bib16,bib17,bib18,bib19,bib20,bib21,bib22}. Although these approaches can produce dense per-frame depth, the connection between front-end tracking and back-end dense recovery remains weak. In particular, structural information estimated by the front end, including local topology, sparse physical anchors, tracking confidence, and tissue-dependent physical variation, is rarely transferred to the back end. As a result, dense predictions may appear visually plausible yet remain misaligned with front-end tracking in affine scale and shift as well as local surface geometry across frames.

In this work, we address monocular endoscopic 3D reconstruction from the perspective of dense and scale-consistent intraoperative point cloud generation. The proposed method builds on a monocular non-rigid SLAM front end and a diffusion-based depth completion back end, but its core contribution is not a simple combination of the two. Instead, the front end estimates tissue-aware heterogeneous physical priors and organizes them into a tissue-aware heterogeneous physical graph (THPG), whose local topology, sparse physical anchors, tracking states, and tissue-dependent physical variation are explicitly propagated to the back end to guide test-time dense depth completion. In this way, sparse deformation estimation and dense depth recovery are coupled through a shared graph representation, with the aim of improving front-end tracking stability and dense point cloud scale consistency under continuous soft-tissue deformation.

The main contributions of this work are summarized as follows:
\begin{enumerate}
    \item We propose a tissue-aware heterogeneous physical graph (THPG) for monocular non-rigid endoscopic tracking. The proposed front end converts coarse tissue-region priors into node-level relative elastic and viscous weights, and integrates these heterogeneous attributes with deformation-consistency-based topology updating. Compared with homogeneous viscoelastic regularization, THPG provides spatially varying physical constraints for improving camera-pose estimation, sparse-anchor reliability, and local topology recovery under heterogeneous soft-tissue deformation.

    \item We develop a THPG-conditioned diffusion depth completion strategy to couple sparse non-rigid tracking with dense pseudo-metric depth recovery. Instead of treating sparse depth samples as isolated guidance, the proposed back end propagates front-end anchors, graph topology, node-level physical attributes, and tracking confidence into DDIM-based test-time optimization. This design is intended to improve the alignment between dense depth prediction and front-end sparse geometry while reducing affine scale-and-shift mismatch and cross-frame local geometric inconsistency without retraining the foundation depth model.

    \item We construct a scale-consistent monocular endoscopic reconstruction method for generating dense and scale-consistent intraoperative point clouds. The complete method integrates THPG-based front-end tracking, physics-guided diffusion depth completion, multi-frame fusion, and an auxiliary anchor-preserving post-fusion refinement step. Experiments on public endoscopic sequences and a self-collected deformable organ phantom dataset report improvements in non-rigid tracking, dense depth completion, and scale-consistent point cloud generation under the evaluated settings.
\end{enumerate}

\section{Related work}\label{sec:related}

Existing studies on monocular endoscopic 3D reconstruction can be reviewed along a progressive technical line: from rigid or weakly deformable reconstruction, to non-rigid SLAM front ends for sparse tracking, and further to dense depth recovery and completion. This organization is consistent with the reconstruction-level motivation of this work: the intraoperative structure must first be tracked robustly, then densified into a surface with sufficient geometric coverage, and finally kept scale-consistent across frames so that it can serve as a reliable point cloud representation for subsequent geometric analysis and surgical-navigation research.

\subsection{From rigid reconstruction to non-rigid endoscopic tracking}

Early endoscopic 3D reconstruction methods were largely developed from conventional multi-view geometry, structure-from-motion, and SLAM techniques. Feature-based SLAM methods estimate camera motion and sparse structure by detecting, matching, and triangulating image features across frames \cite{bib3,bib4}. These methods have achieved strong performance in static or approximately rigid scenes, but their assumptions are difficult to satisfy in monocular endoscopy. Soft tissue surfaces undergo persistent non-rigid motion caused by respiration, heartbeat, and instrument interaction. Such deformation violates the static-scene assumption and may lead to pose drift, unstable sparse maps, and distorted local geometry.

Another line of dynamic reconstruction methods explicitly models non-rigid deformation. Representative systems such as DynamicFusion, VolumeDeform, KillingFusion, and SurfelWarp established important paradigms for deformation representation, online model updating, and surface reconstruction \cite{bib23,bib24,bib25,bib26}. These methods demonstrate that deformation modeling is essential for dynamic 3D reconstruction. However, they generally rely on RGB-D or stereo observations, or require depth measurements that are not readily available in standard monocular endoscopic settings. Therefore, although they provide useful deformation-representation ideas, they cannot be directly transferred to monocular endoscopic reconstruction where absolute depth is unavailable and tissue motion is continuously changing.

To address the limitations of rigid reconstruction in deformable scenes, several monocular non-rigid SLAM methods have been proposed. DefSLAM introduced monocular tracking and mapping for deformable scenes by imposing local geometric constraints, while SD-DefSLAM improved semi-direct tracking for deformable and intracorporeal scenes \cite{bib5,bib6}. DSDT and MCPD further enhanced deformable tracking through sparse direct information, point-based deformation constraints, and dynamic map maintenance \cite{bib7,bib8}. More recently, NR-SLAM investigated a monocular non-rigid SLAM formulation for medical image sequences, in which deformable tracking and mapping are coupled with graph-based viscoelastic regularization \cite{bib9}. These methods significantly improve front-end robustness compared with rigid SLAM and are highly relevant to monocular endoscopic reconstruction.

However, existing non-rigid SLAM methods mainly focus on front-end pose estimation and sparse geometric recovery. The deformation graph, sparse map points, local topology, and physical regularization terms estimated during tracking are usually used only to stabilize the front end. They are rarely propagated further as structured priors for dense reconstruction. More importantly for dense reconstruction, the sparse map produced by the front end often lacks sufficient surface coverage and local neighborhood density to provide a complete surface representation. As a result, even when the front end provides relatively stable sparse tracking, the subsequent dense reconstruction stage may still be weakly constrained by the non-rigid motion state estimated from the same sequence.

\subsection{Monocular dense depth recovery and diffusion-based completion}

Dense depth recovery is required to obtain spatially complete 3D structure from monocular endoscopic videos. From a reconstruction-level perspective, this stage converts sparse tracking map points into an intraoperative surface with denser coverage and more continuous local geometry. Existing learning-based methods include supervised depth estimation, self-supervised depth estimation, geometry-aware depth prediction, and sparse-to-dense depth completion. Supervised methods usually rely on synthetic RGB-D data, transferred supervision, or domain adaptation to overcome the scarcity of accurate depth labels in endoscopic scenes \cite{bib10,bib11,bib12,bib13,bib14,bib15,bib16,bib17,bib18,bib19}. These methods can generate dense per-frame predictions, but their performance is often affected by domain shift between synthetic training data and real surgical videos.

Self-supervised monocular depth methods reduce the dependence on dense ground-truth depth by exploiting photometric or geometric consistency across adjacent frames \cite{bib21,bib22}. In endoscopic scenarios, however, the assumptions behind photometric consistency are often weakened by tissue deformation, viewpoint-dependent appearance changes, and the eye-in-hand illumination configuration. Geometry-aware methods further introduce structural constraints, such as gradient, normal, surface, or cross-frame consistency, to improve anatomical plausibility and reduce local depth distortion. These approaches improve the internal consistency of predicted depth maps, but most of them still operate within the depth-estimation stage itself and do not explicitly inherit the deformation state estimated by a non-rigid SLAM front end.

Sparse-to-dense depth completion provides another route for dense reconstruction by using sparse geometric measurements as anchors. Recent depth completion methods combine sparse depth samples with RGB images to recover dense depth, and generative depth priors have further improved cross-domain prediction ability \cite{bib10,bib11,bib12,bib13,bib14,bib15,bib16,bib17,bib18,bib19,bib27,bib28,bib29,bib30,bib31,bib32,bib33}. In particular, Marigold-DC shows that a pretrained monocular diffusion depth prior can be adapted to sparse-guided depth completion through test-time optimization without retraining the foundation model \cite{bib20}. This strategy is attractive for endoscopic reconstruction because it can exploit a strong pretrained depth prior while avoiding task-specific retraining.

Nevertheless, existing dense depth recovery and completion methods are usually driven by image evidence, sparse depth points, or generic geometric smoothness. They generally do not explicitly incorporate the local deformation topology, node-level physical attributes, or tracking-state confidence produced by a non-rigid SLAM front end. Consequently, the recovered dense depth may be visually plausible at the single-frame level but still inconsistent with the front-end sparse geometry in affine scale and shift as well as cross-frame local surface structure. This limitation is particularly important in monocular endoscopy, where dense reconstruction should not only be visually reasonable but also remain geometrically compatible with the non-rigid tracking result over time. Otherwise, multi-frame fusion can produce a dense point cloud whose local density is improved but whose scale consistency is insufficient for reliable downstream geometric use.

\subsection{Cross-stage prior propagation for sparse-to-dense reconstruction}

The above studies indicate that front-end non-rigid tracking and back-end dense depth recovery have developed largely as two separate technical lines. Front-end methods improve camera tracking, sparse map stability, and local deformation estimation, whereas back-end methods improve dense depth completeness and visual plausibility. However, in most existing approaches, the connection between these two stages remains weak. Sparse map points may be used as depth samples, but the richer information estimated by the front end, including local graph topology, sparse physical anchors, heterogeneous deformation regularization, and tracking reliability, is rarely transmitted explicitly to the back end.

Recent visual foundation models provide additional opportunities for constructing weak structural priors. SAM and SAM 2 can generate region-level masks without dense manual annotation \cite{bib34,bib35}. For endoscopic reconstruction, the value of such masks does not lie in precise medical semantic recognition alone. Instead, they can provide coarse tissue-region cues that help assign relative physical weights to deformation nodes and form tissue-aware heterogeneous priors. In addition, post-fusion point cloud refinement methods, including score-based denoising, can suppress local noise after multi-frame fusion \cite{bib36,bib37}. In this work, such refinement is treated as an auxiliary post-processing step rather than as an independent network contribution. Generic score-based denoising moves points toward high-probability local surface modes estimated from noisy point cloud distributions, but it does not explicitly preserve the sparse-anchor coordinate convention, front-end deformation topology, or tissue-boundary information required by scale-consistent endoscopic reconstruction. Therefore, we only use score-based point updates as a denoising prior and constrain them with THPG-derived anchors, topology, semantic boundaries, and tracking confidence.

In summary, prior work has made substantial progress in rigid and non-rigid endoscopic tracking, monocular dense depth recovery, sparse-guided depth completion, tissue-region prior extraction, and post-fusion refinement. However, existing methods still lack an explicit mechanism for transmitting front-end non-rigid physical priors to the back-end dense reconstruction stage. This missing cross-stage coupling directly motivates the present work. Instead of treating front-end tracking and back-end depth completion as two weakly related stages, we construct a tissue-aware heterogeneous physical graph in the front end and propagate its topology, sparse physical anchors, node-level physical attributes, and tracking states into diffusion-based depth completion. In this way, the proposed method couples sparse non-rigid estimation with dense pseudo-metric recovery, thereby improving the geometric consistency between front-end tracking and back-end reconstruction under continuous soft-tissue deformation and providing a denser intraoperative point cloud basis for subsequent geometric analysis.

\section{Methods}\label{sec:methods}

\subsection{Problem formulation and method overview}\label{subsec:method_overview}

Following the gap identified above, we formulate monocular endoscopic reconstruction as a cross-stage estimation problem for dense and scale-consistent intraoperative point cloud generation. The front end should provide stable camera poses, sparse geometry, and deformation states under soft-tissue motion, whereas the back end should densify this sparse structure without breaking its pseudo-metric scale convention or local topology. To address the coupled challenges of unstable front-end non-rigid tracking and back-end dense reconstruction that is difficult to align consistently with sparse geometry in affine scale and shift as well as local structure, we propose a unified 3D reconstruction method with tissue-aware heterogeneous physical priors. Given an input video sequence, the method jointly estimates temporally consistent camera poses, local deformation states, and the corresponding dense 3D structure under continuous soft-tissue deformation.

\begin{figure*}[t]
    \centering
    \includegraphics[width=\textwidth]{fig1.pdf}
    \caption{Overview of the proposed method. Given a monocular endoscopic video sequence, the method first extracts tissue-region priors through zero-shot semantic segmentation and maps them to node-level relative physical weights. A tissue-aware heterogeneous physical graph (THPG) is then constructed in the front end, where joint pose--deformation optimization estimates the camera pose, sparse physical anchors, and THPG topology. In the back end, sparse physical anchors, THPG topology, and tracking-state information are incorporated into test-time diffusion-based depth completion to recover dense pseudo-metric depth consistent with front-end motion estimation. Finally, multi-frame fusion is performed to obtain a denser intraoperative point cloud, followed by an auxiliary THPG-guided anchor-preserving refinement step for suppressing residual fusion noise while maintaining anchor consistency and local topology.}
    \label{fig:overview}
\end{figure*}

Methodologically, the proposed method does not simply cascade the non-rigid SLAM backbone of NR-SLAM \cite{bib9} with the test-time diffusion-guidance paradigm of Marigold-DC \cite{bib20}. The key innovation is to convert the front-end deformation representation from a tracking-only regularizer into a transferable tissue-aware heterogeneous physical graph (THPG). In THPG, tissue-region priors, heterogeneous elastic--viscous weights, dynamic topology, high-confidence sparse anchors, and tracking states are organized as a structured physical prior. This prior is explicitly propagated to the back end, where it is used to guide diffusion-based depth completion through sparse anchor alignment, local co-deformation constraints, confidence-aware weighting, and affine scale-and-shift optimization. In this formulation, sparse non-rigid tracking and dense depth recovery are coupled by a shared physical-graph representation rather than treated as two isolated stages. The intended dense point cloud is therefore densified from the sparse SLAM structure while being constrained to preserve its pseudo-metric scale convention and local deformation-aware geometry, which is important for multi-frame fusion and subsequent geometric use.

As shown in Fig.~\ref{fig:overview}, the method first applies zero-shot semantic segmentation to obtain tissue-region priors and maps them to node-level relative physical weights, thereby constructing a tissue-aware heterogeneous physical graph (THPG) in the front end \cite{bib34,bib35}. Joint pose--deformation optimization on this graph then yields stable camera poses, an updated THPG topology, and a set of high-confidence sparse physical anchors. In the back end, these structural cues are incorporated as structured guidance into physics-guided diffusion depth completion. The resulting dense depth maps are subsequently fused across frames to form a denser intraoperative point cloud representation. A THPG-guided anchor-preserving score refinement step is then applied to suppress residual fusion noise while preventing the refinement process from destroying high-confidence sparse anchors or deformation-aware local topology.


\subsection{Rationale of the tissue-aware heterogeneous physical graph}\label{subsec:thpg_rationale}

The motivation for introducing THPG is based on a simple but important property of deformable endoscopic scenes: soft-tissue motion is locally coherent but not globally homogeneous. Neighboring points on the same tissue patch usually exhibit correlated deformation over a short temporal window, whereas points located across anatomical boundaries, tissue folds, or instrument-contact regions may undergo different relative motion. A homogeneous deformation graph with a single global elastic or viscous coefficient imposes the same smoothness strength throughout the scene. This can over-constrain regions where relative motion is physically plausible, while under-constraining stable regions that should remain locally coherent. The resulting optimization may either suppress true local deformation or allow excessive drift, both of which degrade sparse tracking and the subsequent dense reconstruction.

THPG addresses this limitation by converting coarse tissue-region information and tracking reliability into a spatially varying graph regularization prior. In the front end, the graph topology specifies which sparse nodes are expected to co-deform, while the heterogeneous elastic and viscous weights specify how strongly their relative displacement should be preserved. Therefore, the graph does not attempt to estimate true biomechanical parameters of the organ. Instead, it provides a relative physical weighting mechanism: stronger constraints are assigned to locally stable and coherent regions, whereas weaker or dynamically removed constraints are used near boundaries and locally inconsistent motion. From an optimization perspective, THPG acts as a non-uniform graph-Laplacian prior on the deformation field, reducing the degrees of freedom in reliable regions while preserving necessary flexibility in heterogeneous regions.

This design is expected to affect reconstruction through two coupled mechanisms. First, in the front end, heterogeneous regularization reduces the ambiguity between camera motion and tissue deformation by discouraging the optimizer from explaining all reprojection residuals with either a globally smooth deformation field or unstable local node motion. This mechanism is intended to produce more reliable camera poses, sparse anchors, and graph topology. Second, in the back end, the same THPG is propagated as structured guidance for diffusion depth completion. Compared with using sparse depth points alone, THPG provides pairwise local-geometry constraints, tissue-dependent smoothing/restoration strengths, and confidence-aware anchor selection. These constraints are intended to pull the affine-ambiguous diffusion prediction toward the pseudo-metric coordinate convention of the front-end sparse anchors while preserving local surface relationships. Consequently, the dense point cloud is expected to become denser and more consistent with the sparse tracking backbone, which is important for stable multi-frame fusion and later geometric analysis.

\subsection{Tissue-aware heterogeneous physical-graph front end}\label{subsec:frontend}

The front end is designed to support camera pose estimation and sparse geometric recovery in complex non-rigid soft-tissue scenes. Unlike conventional rigid SLAM, the proposed method does not assume global scene rigidity. Instead, local tissue deformation is represented by a tissue-aware heterogeneous physical graph (THPG). Let the set of visible graph nodes at time $t$ be $\mathcal{V}_t$, the edge set be $\mathcal{E}_t$, and the 3D position of node $i$ at time $t$ be $\mathbf{x}_i^t$, with initial reference position $\mathbf{x}_i^0$. The camera projection function is denoted by $\pi(\cdot)$, and the camera back-projection operator is denoted by $\Pi^{-1}(\cdot)$ when dense depth is converted to 3D points. The node displacement is defined as
\begin{equation}
\mathbf{u}_i^t = \mathbf{x}_i^t - \mathbf{x}_i^0.
\label{eq:disp}
\end{equation}

\begin{figure*}[t]
    \centering
    \includegraphics[width=\textwidth]{fig2.pdf}
    \caption{Illustration of the tissue-aware heterogeneous physical-graph front end. The current endoscopic frame is first processed by zero-shot semantic segmentation to obtain tissue-region labels, which are then mapped to node-level relative elastic and viscous weights. Based on these priors, the front end constructs a tissue-aware heterogeneous physical graph (THPG) from sparse map points and establishes graph edges according to local spatial adjacency. A dynamic topology update removes connections that no longer satisfy the co-deformation assumption, and the final optimization jointly estimates camera pose and local deformation by combining observation consistency with heterogeneous viscoelastic regularization.}
    \label{fig:frontend}
\end{figure*}

To obtain tissue-region priors for front-end optimization, we first perform zero-shot semantic segmentation on the input endoscopic image and project the resulting region labels onto THPG nodes. For each visible map node $i$, the projection function $\pi(\cdot)$ maps the node to the image plane, from which the corresponding coarse region label $c_i$ is obtained. A predefined mapping function $\Phi(\cdot)$ is then used to assign a relative stiffness-like weight $k_i$ and a relative damping-like weight $\eta_i$ to each node:
\begin{equation}
(k_i, \eta_i) = \Phi(c_i).
\label{eq:nodeweights}
\end{equation}
where $k_i$ and $\eta_i$ should not be interpreted as measured Young's modulus or true tissue viscosity. Instead, they are normalized stiffness-like and damping-like weights used to modulate graph regularization. A larger $k_i$ imposes stronger resistance to local stretching, whereas a smaller $k_i$ leaves greater flexibility for locally deformable regions. Similarly, a larger $\eta_i$ imposes stronger temporal damping on relative node motion, whereas a smaller $\eta_i$ allows faster local deformation changes. Therefore, $\Phi(\cdot)$ encodes a relative tissue-property tendency rather than absolute biomechanical parameters. This design differs from a homogeneous or globally weighted viscoelastic graph, in which the same elastic and viscous coefficients are applied across the entire scene. By assigning node-level weights before forming edge-level constraints, the proposed front end can better represent local tissue-dependent deformation differences, especially near anatomical boundaries and regions affected by instrument interaction.

Table~\ref{tab:physical_mapping} summarizes the default construction of $\Phi(\cdot)$. The assignment follows a two-level rule. First, the segmentation-derived region label is mapped to a relative tissue-property class according to the expected stiffness or deformability of the local tissue region. Second, boundary, specular, occlusion, instrument-contact, and reprojection-instability cues are used only as reliability corrections: they weaken or remove constraints that are likely to cross different motion regions, but they are not treated as direct measurements of tissue stiffness. The normalized values are kept fixed in all experiments and are used as relative regularization multipliers. Nodes that are unstable under reprojection or temporal consistency checks are excluded from the reliable anchor set instead of introducing an additional node-level confidence variable. Frame-level tracking reliability is not introduced as an additional node-level variable. Instead, it is quantified later by the frame-level confidence $c_t$ in Eq.~(\ref{eq:confidence}) and used to schedule the front-end-dependent back-end loss weights in Eq.~(\ref{eq:scheduler}), while the smoothing and boundary weights are kept fixed as shown in Eq.~(\ref{eq:scheduler-fixed}). This design makes the role of the semantic prior explicit: it provides a weak spatial partition for heterogeneous graph regularization, while the physical weights remain relative optimization weights rather than true material parameters.

\begin{table}[t]
\centering
\caption{Default mapping from tissue-region priors to relative physical weights. The values are normalized regularization multipliers rather than measured biomechanical parameters and are kept fixed across the experiments.}
\label{tab:physical_mapping}
\small
\renewcommand{\arraystretch}{1.15}
\begin{tabular}{p{0.30\textwidth}ccp{0.32\textwidth}}
\toprule
Tissue-region prior or reliability condition & $k_i$ & $\eta_i$ & Interpretation \\
\midrule
Stiffer or locally stable tissue region & 1.00 & 1.00 & Stronger resistance to local stretching and temporally abrupt relative motion. \\
Normally deformable soft-tissue region & 0.75 & 0.75 & Allows moderate deformation while maintaining local motion coherence. \\
Highly deformable tissue fold or contact-affected region & 0.45 & 0.35 & Reduces over-regularization where large local deformation is plausible. \\
Semantic or anatomical boundary neighborhood & 0.30 & 0.25 & Weakens cross-boundary constraints to avoid forcing different tissue parts to co-deform. \\
Specular, occluded, or reprojection-unstable region & not used & not used & Excluded from reliable anchor propagation and strong physical regularization. \\
\bottomrule
\end{tabular}
\end{table}

At the graph level, the THPG serves as the core representation of the non-rigid scene. Graph nodes are formed from the sparse map points estimated in the front end, and graph edges are established according to local 3D adjacency to indicate which nodes should be regularized jointly within the current temporal window. The initial connectivity is determined by spatial proximity, whereas the effective strength of each edge is further modulated by node-level physical attributes. This separates the topological assumption of local co-deformation from the strength of the physical regularization, which is important because geometrically adjacent points may still have different deformation reliability in endoscopic scenes. For an edge connecting nodes $i$ and $j$, the joint elastic $k_{ij}$ and viscous $\eta_{ij}$ weights are defined as
\begin{equation}
k_{ij} = \sqrt{k_i k_j}, \qquad \eta_{ij} = \sqrt{\eta_i \eta_j}.
\label{eq:edgeweights}
\end{equation}
The geometric mean gives a symmetric edge strength and prevents a single node with an extremely large weight from dominating the pairwise constraint. Thus, an edge becomes strongly regularized only when both connected nodes support a coherent local-motion assumption.

As the viewpoint changes and tissue deformation evolves, some graph connections may no longer reflect co-motion on the same local surface and therefore need to be updated dynamically. For an edge $(i,j) \in \mathcal{E}_t$, if the relative change in its current length with respect to the initial length exceeds a threshold $\tau_d$, the edge is considered to violate the co-deformation assumption and is removed:
\begin{equation}
\frac{\left|\left\lVert \mathbf{x}_i^t - \mathbf{x}_j^t \right\rVert - \left\lVert \mathbf{x}_i^0 - \mathbf{x}_j^0 \right\rVert\right|}{\left\lVert \mathbf{x}_i^0 - \mathbf{x}_j^0 \right\rVert} > \tau_d.
\label{eq:topology}
\end{equation}
This dynamic topology update is another reason for the performance gain. It prevents outdated graph edges from forcing two nodes to move together after their relative motion has become inconsistent, thereby reducing error propagation from local deformation to camera-pose estimation and sparse map updating.

At each frame, the front end jointly optimizes an observation-consistency term and heterogeneous viscoelastic regularization terms to estimate the camera pose $T_t \in \mathrm{SE}(3)$ and node-wise local deformation $\{\mathbf{u}_i^t\}$. The overall objective is
\begin{equation}
E_{\mathrm{front}} = E_{\mathrm{obs}} + \lambda_e E_{\mathrm{ela}} + \lambda_v E_{\mathrm{vis}}.
\label{eq:front}
\end{equation}
where $E_{\mathrm{obs}}$ denotes the reprojection term, $E_{\mathrm{ela}}$ the heterogeneous elastic regularization term, $E_{\mathrm{vis}}$ the heterogeneous viscous regularization term, and $\lambda_e$ and $\lambda_v$ their corresponding weights.

The observation-consistency term $E_{\mathrm{obs}}$ enforces agreement between the estimated projection and the corresponding image observation. For the $m$-th matched 3D--2D correspondence, let $\mathbf{X}_m$ be the reference 3D map point and $\mathbf{y}_m$ be its matched 2D image observation. Let $\mathcal{W}_t(\mathbf{X}_m)$ denote the graph-based non-rigid warp at frame $t$, obtained by interpolating the displacements of the neighboring THPG nodes associated with $\mathbf{X}_m$. The warped point is then projected by the current camera pose $T_t$, and the reprojection term is defined as
\begin{equation}
E_{\mathrm{obs}} = \sum_m \rho\!\left(\left(\pi\!\left(T_t\mathcal{W}_t(\mathbf{X}_m)\right) - \mathbf{y}_m\right)^\top \Sigma_m^{-1} \left(\pi\!\left(T_t\mathcal{W}_t(\mathbf{X}_m)\right) - \mathbf{y}_m\right)\right).
\label{eq:obs}
\end{equation}
where $m$ indexes matched 3D--2D correspondences, $\rho(\cdot)$ is the Huber robust kernel, and $\Sigma_m$ is the covariance matrix of the 2D observation. This term is used to reduce the influence of correspondence mismatches and local tracking outliers while explicitly coupling the camera pose and graph-based local deformation in the reprojection residual.

The elastic term $E_{\mathrm{ela}}$ penalizes implausible changes in the static distance between adjacent nodes:
\begin{equation}
E_{\mathrm{ela}} = \sum_{(i,j)\in\mathcal{E}_t} k_{ij} \left\lVert \left(\mathbf{x}_i^t - \mathbf{x}_j^t\right) - \left(\mathbf{x}_i^0 - \mathbf{x}_j^0\right) \right\rVert^2.
\label{eq:ela}
\end{equation}
The viscous term $E_{\mathrm{vis}}$ constrains temporal changes in relative motion between neighboring nodes:
\begin{equation}
E_{\mathrm{vis}} = \sum_{(i,j)\in\mathcal{E}_t} \eta_{ij} \left\lVert \left(\mathbf{u}_i^t - \mathbf{u}_j^t\right) - \left(\mathbf{u}_i^{t-1} - \mathbf{u}_j^{t-1}\right) \right\rVert^2.
\label{eq:vis}
\end{equation}

Unlike NR-SLAM, our method does not rely on uniform global coefficients for elastic and viscous regularization. Instead, both terms are modulated by node-level heterogeneous weights, allowing the deformation model to capture not only geometric adjacency but also tissue-dependent variation. The expected effect is a more balanced bias--variance trade-off in non-rigid optimization: stable regions are constrained against free drift, whereas locally deformable regions are not forced to obey an overly rigid global smoothness prior. After optimization, the front end outputs the current camera pose, the updated THPG topology, and a set of reliable visible graph nodes with consistent reprojection and stable tracking states. These nodes are subsequently used as sparse physical anchors for back-end depth completion.

\subsection{Physics-guided diffusion depth completion back end}\label{subsec:backend}

Although THPG-based front-end tracking provides relatively stable camera poses, local deformation estimates, and sparse geometric anchors, its point density remains insufficient for high-quality dense reconstruction and complete intraoperative surface representation. We therefore introduce a physics-guided diffusion depth completion stage in the back end, which propagates the topology, anchor information, and relative physical attributes estimated by the front end into test-time optimization. The goal is not merely to make each depth map visually dense, but to obtain dense depth that remains aligned with the same sparse-anchor coordinate convention across frames.

\begin{figure*}[t]
    \centering
    \includegraphics[width=\textwidth]{fig3.pdf}
    \caption{Illustration of the physics-guided diffusion depth completion back end. High-confidence sparse physical anchors, THPG topology, node-level physical attributes, and tracking states estimated by the front end are explicitly introduced as structured conditions into the DDIM reverse process. The back end combines geometric consistency, viscous smoothing, elastic restoration, scale alignment, and boundary constraints in a joint objective, thereby improving local geometric stability and pseudo-metric scale alignment in zero-shot dense depth completion.}
    \label{fig:backend}
\end{figure*}

During prior transmission, front-end THPG nodes are not passed to the back end indiscriminately. Instead, we retain only those nodes that are visible in the current frame, consistent under reprojection, and sufficiently reliable in tracking. We denote by $\mathcal{A}_t$ the index set of retained visible anchor nodes at frame $t$. For each anchor node $i\in\mathcal{A}_t$, $d_i$ denotes its pseudo-metric anchor depth, $\mathbf{q}_i$ denotes its projected image coordinate, and $\hat d_i=\hat D_t(\mathbf{q}_i)$ denotes the predicted depth sampled at that coordinate. For a pixel or projected coordinate $p$, the corresponding back-projected 3D point is denoted by $\hat{\mathbf{P}}_p=\Pi^{-1}(p,\hat d_p)$, where $\Pi^{-1}$ uses the camera intrinsics of the current view. Thus, the predicted 3D point corresponding to anchor $i$ is $\hat{\mathbf{P}}_i:=\hat{\mathbf{P}}_{\mathbf{q}_i}$. The active anchor-edge set is defined as $\mathcal{E}_t^{\mathcal{A}}=\{(i,j)\in\mathcal{E}_t\mid i,j\in\mathcal{A}_t\}$, so that graph-guided back-end constraints are applied only to edges whose two endpoints are reliable projected anchors.

Based on these sparse anchors, the back end inherits both local adjacency relations and node-level physical attributes from the front-end THPG, from which it constructs the weights required for geometric consistency, viscous smoothing, and elastic restoration. In particular, the connection weight $w_{ij}$ combines 3D anchor proximity and image-plane proximity so that graph guidance remains robust to both geometric dispersion and projection sparsity. The theoretical advantage over depth-sample-only sparse guidance is that the back end receives not only isolated depth samples but also pairwise structural relations among these samples. Therefore, the completed dense surface is optimized to match both the anchor depths and the local relative geometry encoded by the front-end deformation graph.

Let the depth map predicted by the diffusion model for the current view be $\hat D_t$, its depth value at pixel $p$ be $\hat d_p$, and the corresponding back-projected 3D position be $\hat{\mathbf{P}}_p=\Pi^{-1}(p,\hat d_p)$. The back-end objective is defined as
\begin{equation}
E_{\mathrm{back}} = \lambda_g E_{\mathrm{geo}} + \lambda_s E_{\mathrm{smooth}} + \lambda_r E_{\mathrm{rest}} + \lambda_a E_{\mathrm{align}} + \lambda_b E_{\mathrm{bound}}.
\label{eq:eback}
\end{equation}
where $E_{\mathrm{geo}}$, $E_{\mathrm{smooth}}$, $E_{\mathrm{rest}}$, $E_{\mathrm{align}}$, and $E_{\mathrm{bound}}$ denote the geometric consistency, viscous smoothing, elastic restoration, scale alignment, and boundary terms, respectively, and $\lambda_g$, $\lambda_s$, $\lambda_r$, $\lambda_a$, and $\lambda_b$ are their corresponding weights. Their individual effects are evaluated by the energy-term ablation in Sec.~\ref{subsec:ablation}.

The geometric-consistency term $E_{\mathrm{geo}}$ encourages the local 3D structure back-projected from dense depth to preserve the adjacency relations encoded by the THPG, thereby suppressing local structural discontinuities:
\begin{equation}
E_{\mathrm{geo}} = \sum_{(i,j)\in\mathcal{E}_t^{\mathcal{A}}} w_{ij} \left\lVert \hat{\mathbf{P}}_i - \hat{\mathbf{P}}_j - \left(\mathbf{x}_i^t - \mathbf{x}_j^t\right) \right\rVert^2.
\label{eq:geo}
\end{equation}
where
\begin{equation}
 w_{ij} = \frac{1}{2}\left[\exp\!\left(-\frac{\lVert \mathbf{x}_i^t-\mathbf{x}_j^t \rVert^2}{2\sigma^2}\right) + \exp\!\left(-\frac{\lVert \mathbf{q}_i-\mathbf{q}_j \rVert_2}{\gamma}\right)\right],
\label{eq:wij}
\end{equation}
with $\sigma$ controlling 3D-distance decay and $\gamma$ controlling image-plane distance decay over the projected anchor coordinates $\mathbf{q}_i$ and $\mathbf{q}_j$. This definition follows the implementation used to construct connection weights from 3D anchor distance and 2D image-plane distance.

The viscous-smoothing term $E_{\mathrm{smooth}}$ uses propagated node-level viscous weights to regularize the depth-gradient field and suppress high-frequency local fluctuations in the completed depth map:
\begin{equation}
E_{\mathrm{smooth}} = \sum_{p\in\Omega} \eta(p) \left\lVert \nabla \hat d_p \right\rVert^2.
\label{eq:smooth}
\end{equation}
where $\Omega$ denotes the image domain, $\nabla$ denotes the image-domain spatial gradient operator, and $\eta(p)$ is the local smoothing weight derived from front-end viscous attributes. In implementation, node-level or edge-level viscous attributes are first assigned to the projected anchors and then interpolated over valid pixels using distance-weighted interpolation; pixels without nearby reliable anchors use the default neutral weight. The same notation $\eta(p)$ is therefore a pixel-domain regularization map rather than an additional learned variable.

The elastic-restoration term $E_{\mathrm{rest}}$ encourages the solution in relatively stable regions to return toward a plausible equilibrium state, thereby reducing local perturbations induced by stochastic diffusion sampling:
\begin{equation}
E_{\mathrm{rest}} = \sum_{p\in\Omega} \tilde{k}(p) \left(\hat d_p - d_p^0\right)^2.
\label{eq:rest}
\end{equation}
where $d_p^0$ denotes the prior depth map initialized from a densified sparse-depth seed. The effective restoration weight is defined as $\tilde{k}(p)=k(p)r_{\mathrm{ssim}}(p)$, where $k(p)$ is propagated from node-level elastic attributes using the same distance-weighted interpolation strategy and $r_{\mathrm{ssim}}(p)\in[0,1]$ is a normalized SSIM-based appearance-consistency score. Thus, restoration is stronger in locally reliable regions and weaker in visually inconsistent regions.

The scale-alignment term $E_{\mathrm{align}}$ progressively pulls the affine-invariant diffusion prediction into the pseudo-metric space defined by the sparse anchors:
\begin{equation}
E_{\mathrm{align}} = \sum_{i\in\mathcal{A}_t} \rho\!\left(s_t \hat d_i + b_t - d_i\right).
\label{eq:align}
\end{equation}
where $s_t$ and $b_t$ denote the scale and shift parameters for frame $t$, $d_i$ is the pseudo-metric depth provided by the front-end anchors, and $\rho(\cdot)$ is a robust kernel. During DDIM reverse diffusion, $s_t$ and $b_t$ are updated at each denoising step together with the latent variable; the denoising-step superscript is omitted for notational simplicity. This term provides the main mechanism by which the back end remains aligned with the front end. For subsequent multi-frame fusion and geometric analysis, this alignment is essential: dense maps with different per-frame scales may increase point density, but they do not form a reliable fused surface if their coordinates are inconsistent with the sparse tracking backbone.

The boundary term $E_{\mathrm{bound}}$ constrains the predicted depth to a reasonable numerical range and prevents abnormal values during denoising:
\begin{equation}
E_{\mathrm{bound}} = \sum_{p\in\Omega} \left[\max\left(0, d_{\min} - \hat d_p\right)^2 + \max\left(0, \hat d_p - d_{\max}\right)^2\right].
\label{eq:bound}
\end{equation}
where $d_{\min}$ and $d_{\max}$ are determined from the valid sparse-anchor depth range and the initialized depth range, and are kept fixed during the current denoising process. In this way, the back end preserves the strong visual prior of the diffusion model while reducing scale mismatch and local geometric distortion under the joint effect of graph-structure and pseudo-metric constraints.

Because the pretrained diffusion-based monocular depth model produces relative or affine-invariant depth, explicit scale and shift modeling is required to align its output with front-end sparse anchors. In this paper, pseudo-metric alignment means that dense depth is brought into the same affine scale-and-shift convention defined by the current reliable sparse anchors, rather than being used as an independent relative depth map for each frame. Let the preview depth generated at the current denoising step be $\hat D_t$. Its pseudo-metric form is
\begin{equation}
\hat D_t^{\mathrm{pm}} = s_t \hat D_t + b_t.
\label{eq:pmetric}
\end{equation}
where $s_t$ and $b_t$ are optimized jointly with the depth latent; their denoising-step index is omitted to keep the notation compact. This turns scale alignment from a post hoc correction into progressive guidance and is intended to reduce the risk that densification introduces a frame-dependent scale mismatch that would later degrade multi-frame fusion or downstream geometric use.

For initialization, we adopt min--max alignment based on the sparse-anchor depth range, rather than relying solely on a single least-squares fit. This choice is consistent with Marigold-DC, where min--max initialization is generally more robust when the number of conditioned points is limited or the initial geometry is inaccurate.

During inference, the diffusion backbone itself remains unchanged. Instead, during DDIM reverse diffusion, the latent variable, scale parameter, and shift parameter are jointly optimized so that sparse physical anchors and graph-structure constraints from the front end are progressively injected into the denoising trajectory. Let the latent variable at the $l$-th denoising step be $\mathbf{z}^l$. Its update is
\begin{equation}
\mathbf{z}^{l+1} = \mathbf{z}^l - \alpha_l \nabla_{\mathbf{z}} E_{\mathrm{back}}.
\label{eq:latent}
\end{equation}
where $\alpha_l$ denotes the guidance strength at the current step. The objective $E_{\mathrm{back}}$ is evaluated on the dense depth decoded from the current latent variable $\mathbf{z}^l$ together with the current scale and shift estimates. The scale and shift parameters are optimized under the same objective. This design allows the model to retain the strong visual prior of diffusion while constraining structural flattening and local drift under geometric and scale constraints.

Because front-end tracking quality may vary with deformation intensity and correspondence reliability, we further introduce a dynamic weight-scheduling strategy based on tracking state. Let $N_t$ denote the number of inlier feature correspondences that remain valid after reprojection-consistency filtering in the current frame and let $N_{\mathrm{ref}}$ be the reference threshold. The tracking confidence is defined as
\begin{equation}
c_t = \min\left(\frac{N_t}{N_{\mathrm{ref}}}, 1\right).
\label{eq:confidence}
\end{equation}
In our implementation, the reference threshold was set to $N_{\mathrm{ref}}=200$, following the inlier-count setting adopted by NR-SLAM~\cite{bib9}. This value is not used as a hard tracking-failure threshold. Instead, it serves as a saturation reference for normalizing the frame-level tracking confidence. When $N_t\geq N_{\mathrm{ref}}$, the current frame is regarded as providing sufficiently rich and reliable correspondences, and $c_t$ is saturated to 1. When $N_t<N_{\mathrm{ref}}$, the confidence decreases proportionally, thereby weakening the physical-graph guidance in the back-end optimization and preventing unreliable sparse constraints from being over-amplified.

According to the magnitude of $c_t$, the method adaptively adjusts the loss weights that directly depend on front-end tracking reliability. Let $\lambda_g^0$, $\lambda_s^0$, $\lambda_r^0$, $\lambda_a^0$, and $\lambda_b^0$ denote the base weights of the geometric-consistency, viscous-smoothing, elastic-restoration, anchor-alignment, and depth-boundary terms, respectively. In implementation, only the terms that are directly driven by front-end THPG topology, sparse anchors, and physical attributes are scheduled by the frame-level tracking confidence:
\begin{equation}
\lambda_g(t)=c_t\lambda_g^0,\qquad
\lambda_r(t)=c_t\lambda_r^0,\qquad
\lambda_a(t)=c_t\lambda_a^0.
\label{eq:scheduler}
\end{equation}
The smoothing and boundary weights are kept fixed:
\begin{equation}
\lambda_s(t)=\lambda_s^0,\qquad
\lambda_b(t)=\lambda_b^0.
\label{eq:scheduler-fixed}
\end{equation}
Here, $\lambda_s$ remains active to suppress high-frequency depth fluctuations even when the sparse front-end guidance is weak, whereas $\lambda_b$ serves as a numerical range constraint and is therefore not relaxed by tracking confidence. Thus, all five terms in Eq.~(\ref{eq:eback}) participate in back-end optimization, but only the graph-, restoration-, and anchor-related terms are confidence-scheduled. For notational simplicity, Eq.~(\ref{eq:eback}) omits the explicit time superscript of these scheduled weights. When front-end tracking is stable, anchors are sufficiently dense, and reprojection error is low, the geometric-consistency, elastic-restoration, and scale-alignment terms are strengthened so that the back end is driven more strongly by the front-end physical-graph prior. When tracking confidence is low, the influence of these front-end-dependent constraints is relaxed and greater freedom is left to the image prior, preventing sparse but unreliable conditions from being over-amplified. In this way, the back end can exploit informative geometric cues when the front end is reliable, while remaining robust when the front end is uncertain.

\subsection{Auxiliary THPG-guided anchor-preserving point cloud refinement}\label{subsec:postprocess}

After dense depth completion, multi-frame back-projection, and fusion, the reconstructed point cloud may still contain residual high-frequency noise introduced by depth-completion uncertainty, accumulated fusion error, and imperfect non-rigid deformation estimation. A direct application of generic point cloud denoising can improve visual smoothness, but it may also move high-confidence sparse anchors, alter local edge relationships, or slightly change the pseudo-metric scale convention established by the front end. Such changes are undesirable for scale-consistent reconstruction, because the fused point cloud must remain geometrically compatible with the sparse tracking backbone rather than merely appear smoother.

We therefore use score-based point cloud denoising as an auxiliary THPG-guided anchor-preserving refinement step. The score estimator follows the existing score-based denoising principle: a noisy point cloud is treated as samples from a noise-convolved distribution, and each point is moved along the estimated score direction toward a local surface mode \cite{bib36}. This pretrained score estimator is not claimed as a new network or a standalone denoising contribution in this paper. Its role is limited to post-fusion noise suppression. The proposed adaptation constrains the score-based point update with front-end physical anchors, THPG topology, semantic-boundary information, and tracking confidence, so that the optional post-fusion refinement does not destroy the scale convention and local topology established by the front end and the depth-completion back end.

Let \(\mathcal{P}^{0}=\{\mathbf{p}^{0}_{i}\}_{i=1}^{N}\) denote the fused point cloud before refinement, and let \(\mathbf{p}_i\) denote the current optimized position of point \(i\) during refinement, initialized by \(\mathbf{p}_i^{0}\). Let \(\mathcal{A}\) denote the set of fused points associated with reliable sparse physical anchors after multi-frame fusion. For \(i\in\mathcal{A}\), \(\mathbf{a}_{i}\) denotes the corresponding anchor reference position. To avoid confusion with the semantic label \(c_i\) used in Eq.~(\ref{eq:nodeweights}), we use \(\chi_i\in[0,1]\) to denote the local anchor or fusion reliability of point \(i\). The anchor-preserving term $E_{\mathrm{anc}}$ is defined as
\begin{equation}
E_{\mathrm{anc}} = \sum_{i\in\mathcal{A}} \chi_i \left\lVert \mathbf{p}_{i}-\mathbf{a}_{i} \right\rVert_2^2.
\label{eq:refine-anchor}
\end{equation}
This term prevents the refinement step from pulling high-confidence sparse anchors away from the pseudo-metric positions that have already been used to align dense depth in the back end.

To preserve deformation-aware local structure, we further construct a refinement edge set \(\mathcal{E}_{\mathrm{ref}}\) from the active THPG edges. When a dense THPG edge file is unavailable after projection and fusion, image-domain neighboring valid pixels are used as a lightweight approximation, while edges crossing semantic boundaries or large depth discontinuities are suppressed. The topology-preserving term $E_{\mathrm{top}}$ is
\begin{equation}
E_{\mathrm{top}} = \sum_{(i,j)\in\mathcal{E}_{\mathrm{ref}}} \omega_{ij}
\left(\left\lVert \mathbf{p}_{i}-\mathbf{p}_{j}\right\rVert_2 - \ell_{ij}^{0}\right)^2,
\label{eq:refine-topology}
\end{equation}
where \(\ell_{ij}^{0}=\lVert\mathbf{p}_{i}^{0}-\mathbf{p}_{j}^{0}\rVert_2\) is the reference local edge length before score refinement. The edge weight is defined as
\begin{equation}
\omega_{ij}=m_{ij}\sqrt{\chi_i \chi_j}\exp\!\left(-\frac{\lVert\mathbf{p}_{i}^{0}-\mathbf{p}_{j}^{0}\rVert_2^2}{2\sigma_{p}^{2}}\right),
\label{eq:refine-weight}
\end{equation}
where \(\sigma_p\) is the distance bandwidth of the refinement-edge weight and is set from the median reference edge length in \(\mathcal{E}_{\mathrm{ref}}\). The variable \(m_{ij}\) is a boundary-validity indicator. It is set to one for edges within the same semantic or THPG-connected region and reduced to zero for edges crossing semantic boundaries, dynamically removed THPG edges, or strong depth jumps. This design is intended to reduce boundary erosion around tissue folds and contact regions, which is a common risk when generic smoothing is applied to fused endoscopic point clouds.

At refinement iteration \(r\), let \(\mathbf{p}_{i}^{(r)}\) be the current point position. Given the score \(\hat{\mathbf{s}}_{\theta}(\mathbf{p}_{i}^{(r)})\) estimated by the score-based denoising network, the constrained refinement update is written as
\begin{equation}
\mathbf{p}_{i}^{(r+1)} = \mathbf{p}_{i}^{(r)}
+ \alpha^{(r)}\beta_i\hat{\mathbf{s}}_{\theta}\!\left(\mathbf{p}_{i}^{(r)}\right)
- \alpha^{(r)}\nabla_{\mathbf{p}_{i}}
\left(\lambda_{\mathrm{anc}}E_{\mathrm{anc}} + \lambda_{\mathrm{top}}E_{\mathrm{top}}\right),
\label{eq:refine-update}
\end{equation}
where \(\alpha^{(r)}\) is a decreasing step size, \(\lambda_{\mathrm{anc}}\) and \(\lambda_{\mathrm{top}}\) are the anchor- and topology-preservation weights, and \(\beta_i=1-\chi_i\) controls the strength of score-based displacement according to local reliability. The score estimator \(\hat{\mathbf{s}}_{\theta}\) uses fixed pretrained parameters \(\theta\). High-confidence anchors are refined conservatively, whereas low-confidence fused points are allowed stronger score-based correction. Thus, the auxiliary refinement step suppresses local noise while reducing the risk of excessive shrinkage, anchor drift, and topology distortion.

Finally, because the reconstruction is used as a pseudo-metric point cloud rather than an unconstrained shape sample, an optional affine depth correction is applied along the camera ray after refinement:
\begin{equation}
(s^{*},b^{*}) = \arg\min_{s,b}\sum_{i\in\mathcal{A}}\chi_i\left(s z_i+b-z_i^{a}\right)^2,
\label{eq:refine-affine}
\end{equation}
where \(z_i\) is the refined depth of point \(i\) obtained by projecting the refined point along the current camera ray, and \(z_i^{a}\) is the corresponding sparse-anchor depth. The corrected depth is then \(z_i\leftarrow s^{*}z_i+b^{*}\). This step keeps the post-processed point cloud in the same sparse-anchor scale convention as the front-end and back-end reconstruction. In summary, this subsection should be interpreted as a task-specific, consistency-preserving use of score-based point cloud denoising after fusion, rather than as a new generic score-estimation network or a primary contribution of the method.

\subsection{Implementation summary}\label{subsec:implementation}

In implementation, the proposed method follows a two-stage design. For each input frame, a zero-shot segmentation model first produces tissue-region masks, and the visible THPG nodes are projected onto the image plane to inherit tissue-region labels and the corresponding relative elastic and viscous weights through the mapping function $\Phi(\cdot)$. The front-end optimizer then estimates the camera pose $T_t$, local node displacements, the updated graph topology, and the set of reliable visible anchor nodes. These front-end outputs are exported to the back end as sparse anchor depths, active THPG edges, node attributes, and tracking confidence. To build the back-end guidance terms, the connection weights are computed from both 3D anchor distance and 2D image-plane distance according to Eq.~(\ref{eq:wij}); viscous attributes are converted to connection-wise weights over active anchor pairs; and the elastic restoration weights are formed from the elastic coefficient together with an SSIM-based reliability term. The prior depth $d_p^0$ is initialized by densifying the sparse depth with normalized Gaussian convolution, while original sparse measurements are retained at valid anchor pixels. During DDIM inference, the latent depth variable and the affine scale and shift parameters are jointly optimized under Eq.~(\ref{eq:eback}). The completed depth maps are then back-projected and fused across frames. For the optional post-fusion refinement, the pretrained score-based denoiser first provides an initial noise-suppressed point set, after which Eqs.~(\ref{eq:refine-anchor})--(\ref{eq:refine-affine}) are applied to preserve sparse anchors, local topology, semantic boundaries, and the pseudo-metric depth convention.

Table~\ref{tab:notation_params} summarizes the main variables and hyperparameters whose values or construction rules are not self-evident from the equations. All scalar trade-off coefficients are selected once on the validation split and then kept fixed for all test sequences. Data-dependent bandwidths are computed from the active anchors or refinement edges of the current frame, which avoids introducing sequence-specific manual tuning.

% Table~\ref{tab:notation_params} is reserved here but placed at the manuscript end because of its size.
\addtocounter{table}{1}

\section{Experiments}\label{sec:experiments}

\subsection{Experimental setup}\label{subsec:setup}

Experiments were conducted on two public datasets and one self-collected auxiliary dataset to evaluate the proposed method from three perspectives: public benchmark comparability, cross-dataset generalization, and stability under controllable deformation. Since the intended output is a dense and scale-consistent intraoperative point cloud, the evaluation focuses on two reconstruction-level properties: stable sparse tracking and scale-consistent dense recovery. Direct preoperative-to-intraoperative registration is outside the experimental scope of this study and is treated as a potential downstream application. The Hamlyn laparoscopic dataset was used to evaluate front-end non-rigid tracking, with abdominal-sequence-6, exploration-sequence-20, and liver-seq-21 selected for analysis. For back-end zero-shot depth completion, Hamlyn exploration-sequence-20 and SCARED Dataset~1, Sequence~2 were used. In addition, a self-collected deformable organ phantom dataset was constructed to provide auxiliary validation under controllable periodic deformation.

\begin{figure*}[t]
    \centering
    \includegraphics[width=0.95\textwidth]{fig6.pdf}
    \caption{Experimental setup of the self-collected deformable organ phantom platform. A monocular endoscope is inserted into the deformable organ phantom to acquire the endoscopic image stream shown on the display terminal, while the actuation modules generate controllable phantom deformation. For auxiliary pose evaluation, the optical tracking system tracks the four-marker cluster attached to the tail of the endoscope, from which the endoscope pose is estimated through a fixed rigid transformation. The optical tracking system is not used as a depth sensor or as an input to the proposed monocular reconstruction method.}
    \label{fig:experimental-setup}
\end{figure*}

As shown in Fig.~\ref{fig:experimental-setup}, the self-collected platform consists of a deformable organ phantom, a monocular endoscope, actuation modules, a display terminal, and an optical tracking system. During acquisition, the monocular endoscope provides the only image stream used by the reconstruction method, whereas the actuation modules introduce controlled deformation to emulate dynamic organ motion. A four-marker cluster is rigidly attached to the tail of the endoscope and tracked by the optical tracking system. The tracked marker-cluster pose is then transformed to the endoscope coordinate frame using a fixed calibration, providing a reference endoscope pose for auxiliary evaluation. This design allows the self-collected dataset to assess tracking stability under controllable non-rigid deformation, but it does not provide dense reference depth and does not change the monocular input setting of the proposed method.

In the front-end experiments, the proposed method was compared with ORB-SLAM3 \cite{bib4}, SD-DefSLAM \cite{bib6}, MCPD \cite{bib8}, and NR-SLAM \cite{bib9}. For public sequences such as Hamlyn that provide reference depth, reconstruction root mean square error (RMSE) was used to assess geometric recovery accuracy, and the number of successfully tracked frames was also reported. In the self-collected organ phantom experiments, absolute trajectory error (ATE) was computed with respect to the optical-tracking reference poses, and the successful tracking rate was also reported as an auxiliary metric. In the back-end experiments, the method was compared with representative supervised depth completion methods and zero-shot diffusion-based baselines, including Marigold+LS and Marigold-DC \cite{bib13,bib16,bib17,bib20,bib27}. Mean absolute error (MAE) and RMSE were used for quantitative evaluation on datasets with available depth references.

The method was implemented using a hybrid C++/PyTorch codebase. The front-end non-rigid tracking stage runs on the CPU, whereas the back-end diffusion depth completion stage and the auxiliary THPG-guided score refinement step are implemented in PyTorch and accelerated on a GPU. All experiments were performed on Ubuntu 22.04.5 LTS with an Intel Xeon Gold 6248R CPU and an NVIDIA RTX 3090 GPU.

In the back end, a DDIM sampler was used for test-time optimization with 50 inference steps. Following the default setting of Marigold-DC, the learning rate for the latent depth variable was set to 0.05 and that for the scale-and-shift parameters to 0.005. For methods requiring scale recovery, scale fitting was performed uniformly within the common support defined by the predicted depth and sparse anchors to ensure a fair comparison. In the optional final refinement stage, the score-based denoising network was used without retraining. The constraints were applied only during inference, using front-end anchors, projected THPG or valid-pixel topology, semantic boundary masks when available, and confidence-weighted anchor preservation. This stage was not evaluated as an independent denoising method; it was included to keep post-fusion smoothing consistent with the sparse-anchor coordinate convention.

\subsection{Front-end results: non-rigid tracking and sparse geometric recovery}\label{subsec:frontend-results}

This experiment investigates whether, on top of a THPG-based tracking backbone, tissue-aware heterogeneous node weights are associated with lower front-end geometric recovery errors under complex soft-tissue boundaries and local relative motion. To this end, the proposed method was compared with ORB-SLAM3 \cite{bib4}, SD-DefSLAM \cite{bib6}, MCPD \cite{bib8}, and NR-SLAM \cite{bib9} on the Hamlyn dataset.

\begin{table*}[t]
\caption{Comparison of front-end non-rigid tracking performance on the Hamlyn dataset.}
\label{tab:frontend}
\centering
\small
\resizebox{\textwidth}{!}{%
\begin{tabular}{lcccccc}
\toprule
Method & \makecell{abdominal-seq-6\\RMSE$\downarrow$} & \makecell{abdominal-seq-6\\\#Fr.$\uparrow$} & \makecell{exploration-seq-20\\RMSE$\downarrow$} & \makecell{exploration-seq-20\\\#Fr.$\uparrow$} & \makecell{liver-seq-21\\RMSE$\downarrow$} & \makecell{liver-seq-21\\\#Fr.$\uparrow$} \\
\midrule
ORB-SLAM3~\cite{bib4} & 4.85 & 128 & 1.37 & 220 & -- & -- \\
SD-DefSLAM~\cite{bib6} & 2.72 & 286 & 4.68 & 252 & 6.19 & 323 \\
MCPD~\cite{bib8} & -- & -- & 1.48 & 350 & 1.55 & 300 \\
NR-SLAM~\cite{bib9} & 2.37 & 1236 & 1.09 & 609 & 0.59 & 376 \\
Ours & \textbf{1.66} & \textbf{1236} & \textbf{0.76} & \textbf{609} & \textbf{0.42} & \textbf{376} \\
\botrule
\end{tabular}}
\end{table*}

Table~\ref{tab:frontend} reports the front-end comparison on three Hamlyn sequences. ORB-SLAM3 maintains tracking on segments with relatively mild deformation and smooth viewpoint variation, but it degrades or fails in the tested sequences with stronger non-rigid motion. SD-DefSLAM reduces the error on abdominal-sequence-6 but remains less stable on exploration-sequence-20 and liver-seq-21. MCPD performs competitively on the dynamic sequences where results are available, while NR-SLAM provides a stronger non-rigid SLAM baseline by combining graph-based deformation representation with viscoelastic regularization.

Relative to NR-SLAM, the proposed method reports lower RMSE on the three tested public sequences: from 2.37~mm to 1.66~mm on abdominal-sequence-6, from 1.09~mm to 0.76~mm on exploration-sequence-20, and from 0.59~mm to 0.42~mm on liver-seq-21. These results suggest that tissue-aware heterogeneous node weights are beneficial for local structural stability and geometric consistency in the evaluated deformation scenarios.

\begin{table*}[t]
\caption{Auxiliary quantitative validation on the self-collected deformable organ phantom dataset.}
\label{tab:phantom}
\centering
\small
\resizebox{\textwidth}{!}{%
\begin{tabular}{lcccc}
\toprule
Method & \makecell{stomach phantom\\ATE (mm)$\downarrow$} & \makecell{stomach phantom\\successful tracking rate (\%)$\uparrow$} & \makecell{cardiac phantom\\ATE (mm)$\downarrow$} & \makecell{cardiac phantom\\successful tracking rate (\%)$\uparrow$} \\
\midrule
ORB-SLAM3~\cite{bib4} & 3.42 & 82.6 & 4.18 & 74.3 \\
SD-DefSLAM~\cite{bib6} & 2.53 & 91.6 & 3.07 & 88.9 \\
MCPD~\cite{bib8} & 2.06 & 95.1 & 2.58 & 93.4 \\
NR-SLAM~\cite{bib9} & 1.74 & 97.9 & 2.21 & 96.5 \\
Ours & \textbf{1.28} & \textbf{99.1} & \textbf{1.63} & \textbf{98.6} \\
\botrule
\end{tabular}}
\end{table*}

The self-collected phantom data in Table~\ref{tab:phantom} show a similar ranking under controllable periodic deformation. ORB-SLAM3 exhibits noticeable drift in this setting. SD-DefSLAM and MCPD reduce ATE compared with the rigid baseline, indicating the benefit of explicit deformation modeling. NR-SLAM further reduces the reported ATE, while the proposed method reports the lowest ATE and highest successful tracking rate on both the stomach-phantom and cardiac-phantom sequences within this auxiliary validation.

To further visualize front-end tracking behavior, Fig.~\ref{fig:qual-frontend} presents representative qualitative comparisons on both public endoscopic sequences and the self-collected deformable organ phantom data. Only several representative non-rigid baselines are shown. In the displayed examples, some competing methods show local drift, structural stretching, or scattered point distributions under large deformation, boundary transition, and pronounced viewpoint change, whereas the proposed method produces more compact sparse geometric recovery.

\begin{figure*}[t]
    \centering
    \includegraphics[width=\textwidth]{fig4.pdf}
    \caption{Qualitative comparison of front-end non-rigid tracking results. Representative examples from public endoscopic sequences and the self-collected deformable organ phantom data are shown. The compared methods exhibit different behaviors under complex deformation, boundary transition, and viewpoint change. In the displayed cases, some baselines show local drift, structural stretching, and scattered point distributions, whereas the proposed method maintains more stable local geometry and reconstruction continuity.}
    \label{fig:qual-frontend}
\end{figure*}

\subsection{Back-end results: dense depth completion and scale-aligned recovery}\label{subsec:backend-results}

This experiment examines whether, without retraining the foundation diffusion depth model, explicitly propagating the structured physical priors estimated by the front-end THPG into test-time optimization is associated with lower dense depth errors than generative methods guided only by sparse depth points. This question is directly related to reliable multi-frame reconstruction, because densification is useful only when the completed depth remains compatible with the sparse tracking backbone and does not introduce additional frame-dependent scale mismatch. To answer this question, the proposed method was evaluated on Hamlyn exploration-sequence-20 and SCARED Dataset~1, Sequence~2, and compared with representative supervised depth completion methods as well as zero-shot diffusion-based baselines.

\begin{table*}[t]
\caption{Cross-dataset comparison of dense depth completion errors.}
\label{tab:backend}
\centering
\small
\resizebox{\textwidth}{!}{%
\begin{tabular}{lcccc}
\toprule
Method & \makecell{Hamlyn exploration-sequence-20\\MAE$\downarrow$} & \makecell{Hamlyn exploration-sequence-20\\RMSE$\downarrow$} & \makecell{SCARED Dataset~1, Sequence~2\\MAE$\downarrow$} & \makecell{SCARED Dataset~1, Sequence~2\\RMSE$\downarrow$} \\
\midrule
NLSPN~\cite{bib13} & 3.28 & 5.76 & 4.15 & 6.92 \\
CompletionFormer~\cite{bib16} & 2.87 & 5.12 & 3.64 & 6.35 \\
VPP4DC~\cite{bib17} & 1.86 & 3.59 & 2.27 & 4.18 \\
Marigold+LS~\cite{bib27} & 2.54 & 4.21 & 3.02 & 5.06 \\
Marigold-DC~\cite{bib20} & 1.52 & 2.94 & 1.89 & 3.47 \\
Ours & \textbf{1.18} & \textbf{2.31} & \textbf{1.45} & \textbf{2.76} \\
\botrule
\end{tabular}}
\end{table*}

Table~\ref{tab:backend} reports the dense depth comparison on Hamlyn exploration-sequence-20 and SCARED Dataset~1, Sequence~2. The supervised depth completion baselines yield higher errors in these endoscopic test settings than the zero-shot diffusion-based baselines. Marigold+LS, which combines a pretrained monocular depth prior with a post-hoc least-squares scale fit, recovers a globally plausible coarse depth layout in many cases, but its error remains higher than Marigold-DC and the proposed method under the tested sparse-anchor conditions.

Building on Marigold-DC, the proposed method further incorporates local geometric consistency and heterogeneous physical priors derived from the front-end THPG, extending the constraint from sparse-point fitting to generation-time optimization jointly governed by local graph structure and physical attributes. On Hamlyn exploration-sequence-20, the RMSE decreases from 2.94 to 2.31 relative to Marigold-DC; on SCARED Dataset~1, Sequence~2, it decreases from 3.47 to 2.76. These results indicate that propagating front-end structural priors into the back end is associated with lower dense depth errors and better pseudo-metric scale alignment in the evaluated dynamic monocular endoscopic sequences.

To further compare depth-recovery behavior under extremely sparse-anchor conditions, Fig.~\ref{fig:qual-backend} presents representative qualitative examples. In the displayed cases, several supervised depth completion methods show blurred boundaries, over-smoothed local regions, or confused depth ordering in medical images. Some methods recover the overall depth trend, but local scale instability or geometric collapse remains visible when anchors are sparse, local deformation changes rapidly, or boundary geometry becomes ambiguous. The proposed method shows improved organ-boundary preservation, local surface undulation, and overall depth continuity in these examples.

% Fig.~\ref{fig:qual-backend} is placed at the manuscript end because of its large visual footprint.

\subsection{Cross-stage benefit analysis of the complete method}\label{subsec:system}

After the separate front-end and back-end comparisons, we further analyze the interaction between the two stages within the complete method. To this end, three representative variants were compared: (1) front-end sparse non-rigid tracking only; (2) front-end + vanilla diffusion completion; and (3) front-end + physics-guided diffusion completion.

\begin{table}[t]
\caption{Cross-stage benefit analysis on representative sequences.}
\label{tab:system}
\centering
\small
\begin{tabular}{p{0.42\columnwidth}cc}
\toprule
Method variant & \makecell{Hamlyn liver-seq-21\\front-end RMSE (mm)$\downarrow$} & \makecell{SCARED seq2\\depth RMSE (mm)$\downarrow$} \\
\midrule
Front-end only & 0.42 & -- \\
Front-end + vanilla diffusion & 0.42 & 3.18 \\
Front-end + physics-guided diffusion & 0.42 & 2.76 \\
\botrule
\end{tabular}
\end{table}

The comparison in Table~\ref{tab:system} indicates that the front-end-only configuration can maintain basic pose estimation and local geometric constraints, but it does not provide dense depth recovery. Introducing vanilla diffusion completion on top of the front end reduces the lack of dense surface coverage, but the depth RMSE remains 3.18 on SCARED seq2 when no front-end structural prior is injected. The physics-guided diffusion variant reduces this value to 2.76 in the same setting, supporting the value of explicit cross-stage prior propagation in the evaluated method.


\subsection{Ablation study}\label{subsec:ablation}

To more explicitly examine the contribution of the proposed components, we expanded the ablation analysis into three levels. First, a front-end ablation was conducted to isolate the effects of dynamic topology updating and heterogeneous elastic--viscous weighting. Second, a back-end prior-propagation ablation was conducted to separate sparse-anchor alignment, THPG topology, node-level physical attributes, and confidence-aware scheduling. Third, because the back-end objective in Eq.~(\ref{eq:eback}) contains five energy terms, an energy-term ablation was conducted by removing one term at a time while keeping the remaining settings unchanged. The post-fusion refinement step is not treated as a separate denoising benchmark in this ablation study, because it uses a pretrained score estimator and only adds anchor- and topology-preserving constraints during inference. Therefore, the following ablations focus on the front-end and back-end components that form the main methodological contribution, while the refinement step is kept as a fixed auxiliary consistency-preserving post-processing step.

\begin{table}[t]
\caption{Fine-grained front-end ablation on Hamlyn liver-seq-21. Each variant changes only one part of the THPG front end.}
\label{tab:frontend_ablation}
\centering
\small
\renewcommand{\arraystretch}{1.12}
\begin{tabular}{p{0.62\columnwidth}c}
\toprule
Variant & \makecell{RMSE (mm)$\downarrow$} \\
\midrule
NR-SLAM baseline~\cite{bib9} & 0.59 \\
Homogeneous physical graph & 0.56 \\
w/o dynamic topology update & 0.50 \\
w/o heterogeneous elastic weight & 0.48 \\
w/o heterogeneous viscous weight & 0.46 \\
Full THPG front end & \textbf{0.42} \\
\botrule
\end{tabular}
\end{table}

Table~\ref{tab:frontend_ablation} reports the contribution of the proposed THPG front end. The NR-SLAM baseline uses a homogeneous viscoelastic regularization strategy and obtains an RMSE of 0.59~mm on liver-seq-21. Replacing it with a homogeneous physical graph reduces the error to 0.56~mm, suggesting that graph construction alone provides only a moderate gain. Removing the dynamic topology update increases the error to 0.50~mm, which is consistent with the interpretation that outdated edges may force nodes with inconsistent relative motion to co-deform. Removing either the heterogeneous elastic or viscous weight also increases the error relative to the full THPG front end, with a larger degradation observed when the heterogeneous elastic weight is removed. This suggests that stiffness-like spatial regularization plays a more direct role in preserving local geometry, whereas damping-like temporal regularization further stabilizes node motion across frames. The full THPG front end yields the lowest error in this ablation.

\begin{table*}[t]
\caption{Fine-grained back-end prior-propagation ablation on representative dense depth completion sequences. The variants progressively add the front-end priors propagated to the diffusion-based depth completion stage.}
\label{tab:backend_prior_ablation}
\centering
\small
\setlength{\tabcolsep}{4pt}
\renewcommand{\arraystretch}{1.12}
\resizebox{\textwidth}{!}{%
\begin{tabular}{lcccc}
\toprule
Variant & \makecell{Hamlyn exp.-seq-20\\MAE$\downarrow$} & \makecell{Hamlyn exp.-seq-20\\RMSE$\downarrow$} & \makecell{SCARED seq2\\MAE$\downarrow$} & \makecell{SCARED seq2\\RMSE$\downarrow$} \\
\midrule
Sparse-anchor alignment only & 1.41 & 2.72 & 1.75 & 3.23 \\
+ THPG topology & 1.32 & 2.55 & 1.63 & 3.05 \\
+ node physical attributes & 1.25 & 2.43 & 1.54 & 2.91 \\
+ confidence-aware scheduling & \textbf{1.18} & \textbf{2.31} & \textbf{1.45} & \textbf{2.76} \\
\botrule
\end{tabular}}
\end{table*}

Table~\ref{tab:backend_prior_ablation} further separates the contribution of each type of prior transmitted from the front end to the back end. Using only sparse-anchor alignment constrains the affine scale and shift of the diffusion prediction, but it does not encode local co-deformation relationships among anchors. Adding THPG topology reduces the error by introducing pairwise graph constraints, which helps preserve local surface structure beyond isolated depth fitting. Adding node-level physical attributes further reduces the error because the back end can apply different restoration and smoothing strengths according to tissue-dependent deformation tendencies. Finally, confidence-aware scheduling yields the lowest error by strengthening graph-based guidance when the front end is reliable and weakening it when sparse anchors may be unstable. This progressive trend supports the interpretation that the improvement comes from explicit cross-stage prior propagation, rather than from simply adding a diffusion depth completion stage after non-rigid tracking.

\begin{table*}[t]
\caption{Energy-term ablation of the back-end objective on representative dense depth completion sequences. Each variant removes only one term from Eq.~(\ref{eq:eback}) while keeping the other terms and the same DDIM test-time optimization settings unchanged.}
\label{tab:energy_ablation}
\centering
\small
\setlength{\tabcolsep}{4pt}
\renewcommand{\arraystretch}{1.12}
\resizebox{\textwidth}{!}{%
\begin{tabular}{lcccc}
\toprule
Variant & \makecell{Hamlyn exp.-seq-20\\MAE$\downarrow$} & \makecell{Hamlyn exp.-seq-20\\RMSE$\downarrow$} & \makecell{SCARED seq2\\MAE$\downarrow$} & \makecell{SCARED seq2\\RMSE$\downarrow$} \\
\midrule
w/o $E_{\mathrm{geo}}$ & 1.35 & 2.58 & 1.68 & 3.05 \\
w/o $E_{\mathrm{smooth}}$ & 1.27 & 2.45 & 1.56 & 2.91 \\
w/o $E_{\mathrm{rest}}$ & 1.31 & 2.50 & 1.62 & 2.98 \\
w/o $E_{\mathrm{align}}$ & 1.49 & 2.88 & 1.83 & 3.35 \\
w/o $E_{\mathrm{bound}}$ & 1.22 & 2.39 & 1.50 & 2.85 \\
Ours & \textbf{1.18} & \textbf{2.31} & \textbf{1.45} & \textbf{2.76} \\
\botrule
\end{tabular}}
\end{table*}

The energy-term ablation results in Table~\ref{tab:energy_ablation} show that the five terms in Eq.~(\ref{eq:eback}) make non-identical contributions. Removing $E_{\mathrm{align}}$ causes the largest degradation because the diffusion prediction then loses the direct pseudo-metric alignment to the front-end sparse anchors. Removing $E_{\mathrm{geo}}$ also increases the error, indicating that pairwise THPG topology provides useful local structural constraints beyond isolated sparse-depth fitting. Removing $E_{\mathrm{rest}}$ weakens the tissue-attribute-guided restoration toward the initialized equilibrium depth and therefore increases local depth perturbations. Removing $E_{\mathrm{smooth}}$ produces a smaller but consistent degradation, suggesting that the viscous smoothing term mainly suppresses high-frequency depth fluctuations. The effect of removing $E_{\mathrm{bound}}$ is the mildest, but the error still increases because abnormal depth values are no longer explicitly penalized during denoising. Overall, these results support the adopted design of the complete back-end objective: $E_{\mathrm{align}}$ provides sparse-anchor scale alignment, $E_{\mathrm{geo}}$ preserves local graph structure, $E_{\mathrm{rest}}$ introduces tissue-dependent restoration, $E_{\mathrm{smooth}}$ stabilizes the depth-gradient field, and $E_{\mathrm{bound}}$ maintains numerical depth-range stability.

The refinement step should therefore be interpreted as an auxiliary consistency-preserving post-processing component. Its purpose is to reduce residual fusion noise without intentionally changing high-confidence sparse anchors or local THPG-derived topology, rather than to compete with generic point cloud denoising methods.

\subsection{Qualitative analysis and discussion}\label{subsec:discussion}

Beyond quantitative evaluation, we further compared reconstruction results on both public endoscopic sequences and the self-collected deformable organ phantom dataset to examine local structural recovery, cross-frame consistency, and final surface quality. The qualitative observations are consistent with the quantitative trends and suggest that the reported gains arise from coordinated interaction between front-end tracking and back-end dense reconstruction rather than from either stage alone. This is consistent with the reconstruction-level motivation of the method: a useful intraoperative point cloud should be dense enough for local surface analysis and sufficiently scale-consistent to avoid confusing reconstruction error with anatomical deformation.

The observed THPG-related gains can be interpreted from the perspective of constrained non-rigid estimation. In deformable endoscopic scenes, the solution space is highly under-constrained because camera motion, tissue deformation, and monocular depth scale can compensate for each other. A homogeneous physical graph reduces this ambiguity only by enforcing the same smoothness assumption everywhere, which may be insufficient when different tissue regions exhibit different deformation behavior. THPG introduces a more informative prior by keeping local co-deformation constraints where they are reliable, weakening or removing them where relative motion is inconsistent, and carrying the resulting sparse physical anchors into dense depth optimization. This mechanism provides a plausible explanation for the observed reductions in sparse tracking and dense reconstruction errors.

On public endoscopic sequences, rigid baselines show local drift, structural misalignment, or surface tearing in some tested frames with persistent non-rigid motion. With the THPG, the front end shows more stable local motion representation near tissue boundaries and more continuous cross-frame surface relationships in the selected examples. In the back-end qualitative comparison, vanilla diffusion completion can generate visually plausible dense depth, but local inconsistency remains visible when sparse anchors are limited or tracking quality deteriorates. By explicitly propagating the structured physical priors estimated by the front end, the proposed method shows improved local surface shape, boundary continuity, and cross-frame scale consistency in the displayed cases. These observations suggest that the fused surface may provide a stronger intraoperative geometric representation for potential downstream use; however, direct preoperative-to-intraoperative registration is not evaluated as an independent task in this study.

On the self-collected phantom data, the proposed method shows relatively stable surface continuity and deformation consistency under controllable dynamic conditions. The auxiliary THPG-guided point cloud refinement step is used to suppress local high-frequency noise after fusion. Unlike a direct generic denoiser, it is constrained not to move high-confidence sparse anchors or destroy local THPG or boundary-aware topology, so that the final surface remains compatible with the pseudo-metric coordinate convention established by the front end and the back end. We do not present this step as a separate denoising contribution; it is a supporting post-processing component used to maintain consistency after multi-frame fusion. The present study still has several limitations, including limited experimental scale and statistical power, the lack of truly end-to-end joint optimization between the front end and the back end, and the absence of a complete downstream registration experiment using paired preoperative and intraoperative point clouds. The method also depends on zero-shot tissue-region priors that may still be affected by strong domain shift. These issues warrant further investigation in future work.

\section{Conclusion}\label{sec:conclusion}

This study investigated monocular non-rigid endoscopic 3D reconstruction for dense and scale-consistent intraoperative point cloud generation. The work focuses on two tightly coupled bottlenecks: unstable front-end non-rigid tracking and back-end dense reconstruction whose predicted depth is difficult to align consistently with front-end sparse geometry in affine scale and shift as well as cross-frame local structure. To address these challenges, we proposed a unified reconstruction method with tissue-aware heterogeneous physical priors. Rather than treating front-end tracking and back-end depth recovery as largely independent stages, the proposed method explicitly propagates structured physical priors estimated in the front end into back-end test-time diffusion depth completion, thereby coupling sparse non-rigid motion estimation with dense pseudo-metric depth recovery within one unified reconstruction method.

Experimental results show lower front-end reconstruction errors and lower back-end dense depth errors than the evaluated baselines under the tested settings, without retraining the foundation depth model. Front-end experiments suggest that heterogeneous node weights contribute to geometric recovery under complex tissue boundaries and local relative motion. Back-end experiments indicate that structured THPG guidance is associated with lower errors than sparse-depth-only diffusion guidance on the evaluated sequences. The auxiliary THPG-guided refinement is used after fusion to suppress residual noise while preserving anchor consistency and local topology. In addition, cross-stage and ablation studies support the role of explicit cross-stage prior propagation and confidence-aware adaptive scheduling in the reported performance.

The proposed method should be regarded as a unified method-level attempt for monocular non-rigid endoscopic 3D reconstruction rather than as a direct model of true biomechanical processes. Its main value in the present study is to transform sparse non-rigid tracking results into a denser and more scale-consistent intraoperative surface representation. Several limitations remain, including limited experimental scale and statistical power, the lack of fully end-to-end joint optimization, the dependence of zero-shot tissue-region priors on foundation-model generalization, and the need for direct validation in a complete preoperative-to-intraoperative registration workflow. Future work will focus on larger-scale validation, tighter cross-stage optimization, and more robust integration of semantic, geometric, physical, and downstream task-level constraints.

\begin{thebibliography}{99}

\bibitem{bib1}
Mountney P, Stoyanov D, Yang GZ. Three-dimensional tissue deformation recovery and tracking. \textit{IEEE Signal Processing Magazine}. 2010;27(4):14--24.

\bibitem{bib2}
Ma R, Wang R, Zhang Y, et al. RNN-SLAM: Reconstructing the 3D colon to visualize missing regions during a colonoscopy. \textit{Medical Image Analysis}. 2021;72:102100.

\bibitem{bib3}
Mur-Artal R, Montiel JMM, Tardos JD. ORB-SLAM: A versatile and accurate monocular SLAM system. \textit{IEEE Transactions on Robotics}. 2015;31(5):1147--1163.

\bibitem{bib4}
Campos C, Elvira R, Rodr\'{i}guez JJG, et al. ORB-SLAM3: An accurate open-source library for visual, visual-inertial, and multimap SLAM. \textit{IEEE Transactions on Robotics}. 2021;37(6):1874--1890.

\bibitem{bib5}
Lamarca J, Parashar S, Bartoli A, et al. DefSLAM: Tracking and mapping of deforming scenes from monocular sequences. \textit{IEEE Transactions on Robotics}. 2020;37(1):291--303.

\bibitem{bib6}
G\'{o}mez-Rodr\'{i}guez JJ, Lamarca J, Morlana J, et al. SD-DefSLAM: Semi-direct monocular SLAM for deformable and intracorporeal scenes. In: \textit{2021 IEEE International Conference on Robotics and Automation (ICRA)}. IEEE; 2021. p. 5170--5177.

\bibitem{bib7}
Lamarca J, Rodr\'{i}guez JJG, Tardos JD, et al. Direct and sparse deformable tracking. \textit{IEEE Robotics and Automation Letters}. 2022;7(4):11450--11457.

\bibitem{bib8}
Rodr\'{i}guez JJG, Montiel JMM, Tard\'{o}s JD. Tracking monocular camera pose and deformation for SLAM inside the human body. In: \textit{2022 IEEE/RSJ International Conference on Intelligent Robots and Systems (IROS)}. IEEE; 2022. p. 5278--5285.

\bibitem{bib9}
Rodriguez JJG, Montiel JMM, Tardos JD. NR-SLAM: Nonrigid monocular SLAM. \textit{IEEE Transactions on Robotics}. 2024;40:4252--4264.

\bibitem{bib10}
Ma F, Karaman S. Sparse-to-dense: Depth prediction from sparse depth samples and a single image. In: \textit{2018 IEEE International Conference on Robotics and Automation (ICRA)}. IEEE; 2018. p. 4796--4803.

\bibitem{bib11}
Cheng X, Wang P, Yang R. Depth estimation via affinity learned with convolutional spatial propagation network. In: \textit{Proceedings of the European Conference on Computer Vision (ECCV)}; 2018. p. 103--119.

\bibitem{bib12}
Cheng X, Wang P, Guan C, et al. CSPN++: Learning context and resource aware convolutional spatial propagation networks for depth completion. In: \textit{Proceedings of the AAAI Conference on Artificial Intelligence}. 2020;34(7):10615--10622.

\bibitem{bib13}
Park J, Joo K, Hu Z, et al. Non-local spatial propagation network for depth completion. In: \textit{European Conference on Computer Vision}. Cham: Springer International Publishing; 2020. p. 120--136.

\bibitem{bib14}
Rho K, Ha J, Kim Y. GuideFormer: Transformers for image guided depth completion. In: \textit{Proceedings of the IEEE/CVF Conference on Computer Vision and Pattern Recognition}. 2022. p. 6250--6259.

\bibitem{bib15}
Yu Z, Sheng Z, Zhou Z, et al. Aggregating feature point cloud for depth completion. In: \textit{Proceedings of the IEEE/CVF International Conference on Computer Vision}. 2023. p. 8732--8743.

\bibitem{bib16}
Zhang Y, Guo X, Poggi M, et al. CompletionFormer: Depth completion with convolutions and vision transformers. In: \textit{Proceedings of the IEEE/CVF Conference on Computer Vision and Pattern Recognition}. 2023. p. 18527--18536.

\bibitem{bib17}
Bartolomei L, Poggi M, Conti A, et al. Revisiting depth completion from a stereo matching perspective for cross-domain generalization. In: \textit{2024 International Conference on 3D Vision (3DV)}. IEEE; 2024. p. 1360--1370.

\bibitem{bib18}
Tang J, Tian FP, An B, et al. Bilateral propagation network for depth completion. In: \textit{Proceedings of the IEEE/CVF Conference on Computer Vision and Pattern Recognition}. 2024. p. 9763--9772.

\bibitem{bib19}
Zuo Y, Deng J. OGNI-DC: Robust depth completion with optimization-guided neural iterations. In: \textit{European Conference on Computer Vision}. Cham: Springer Nature Switzerland; 2024. p. 78--95.

\bibitem{bib20}
Viola M, Qu K, Metzger N, et al. Marigold-DC: Zero-shot monocular depth completion with guided diffusion. In: \textit{Proceedings of the IEEE/CVF International Conference on Computer Vision}. 2025. p. 5359--5370.

\bibitem{bib21}
Ranftl R, Lasinger K, Hafner D, et al. Towards robust monocular depth estimation: Mixing datasets for zero-shot cross-dataset transfer. \textit{IEEE Transactions on Pattern Analysis and Machine Intelligence}. 2020;44(3):1623--1637.

\bibitem{bib22}
Yin W, Wang X, Shen C, et al. DiverseDepth: Affine-invariant depth prediction using diverse data. \textit{arXiv preprint arXiv:2002.00569}. 2020.

\bibitem{bib23}
Newcombe RA, Fox D, Seitz SM. DynamicFusion: Reconstruction and tracking of non-rigid scenes in real time. In: \textit{Proceedings of the IEEE Conference on Computer Vision and Pattern Recognition}. 2015. p. 343--352.

\bibitem{bib24}
Innmann M, Zollh\"{o}fer M, Nie\ss ner M, et al. VolumeDeform: Real-time volumetric non-rigid reconstruction. In: \textit{European Conference on Computer Vision}. Cham: Springer International Publishing; 2016. p. 362--379.

\bibitem{bib25}
Slavcheva M, Baust M, Cremers D, et al. KillingFusion: Non-rigid 3D reconstruction without correspondences. In: \textit{Proceedings of the IEEE Conference on Computer Vision and Pattern Recognition}. 2017. p. 1386--1395.

\bibitem{bib26}
Gao W, Tedrake R. SurfelWarp: Efficient non-volumetric single view dynamic reconstruction. \textit{arXiv preprint arXiv:1904.13073}. 2019.

\bibitem{bib27}
Zhao W, Rao Y, Liu Z, et al. Unleashing text-to-image diffusion models for visual perception. In: \textit{Proceedings of the IEEE/CVF International Conference on Computer Vision}. 2023. p. 5729--5739.

\bibitem{bib28}
Duan Y, Guo X, Zhu Z. DiffusionDepth: Diffusion denoising approach for monocular depth estimation. In: \textit{European Conference on Computer Vision}. Cham: Springer Nature Switzerland; 2024. p. 432--449.

\bibitem{bib29}
Saxena S, Kar A, Norouzi M, et al. Monocular depth estimation using diffusion models. \textit{arXiv preprint arXiv:2302.14816}. 2023.

\bibitem{bib30}
Saxena S, Herrmann C, Hur J, et al. The surprising effectiveness of diffusion models for optical flow and monocular depth estimation. \textit{Advances in Neural Information Processing Systems}. 2023;36:39443--39469.

\bibitem{bib31}
Ke B, Obukhov A, Huang S, et al. Repurposing diffusion-based image generators for monocular depth estimation. In: \textit{Proceedings of the IEEE/CVF Conference on Computer Vision and Pattern Recognition}. 2024. p. 9492--9502.

\bibitem{bib32}
Yang L, Kang B, Huang Z, et al. Depth Anything: Unleashing the power of large-scale unlabeled data. In: \textit{Proceedings of the IEEE/CVF Conference on Computer Vision and Pattern Recognition}. 2024. p. 10371--10381.

\bibitem{bib33}
Bochkovskii A, Delaunoy AG, Germain H, et al. Depth Pro: Sharp monocular metric depth in less than a second. \textit{arXiv preprint arXiv:2410.02073}. 2024.

\bibitem{bib34}
Kirillov A, Mintun E, Ravi N, et al. Segment Anything. In: \textit{Proceedings of the IEEE/CVF International Conference on Computer Vision}. 2023. p. 4015--4026.

\bibitem{bib35}
Ravi N, Gabeur V, Hu YT, et al. SAM~2: Segment anything in images and videos. \textit{arXiv preprint arXiv:2408.00714}. 2024.

\bibitem{bib36}
Luo S, Hu W. Score-based point cloud denoising. In: \textit{Proceedings of the IEEE/CVF International Conference on Computer Vision}. 2021. p. 4583--4592.

\bibitem{bib37}
Wang Z, Li M, Xin S, et al. Adaptive and iterative point cloud denoising with score-based diffusion model. In: \textit{Computer Graphics Forum}. 2025;44(6):e70149.

\bibitem{bib38}
Song J, Meng C, Ermon S. Denoising diffusion implicit models. \textit{arXiv preprint arXiv:2010.02502}. 2020.

\end{thebibliography}


\clearpage
% Large floats intentionally placed at the manuscript end.
% Table numbering is preserved because Table~\ref{tab:notation_params} is discussed in the Methods section.
\setcounter{table}{1}
\begin{table*}[p]
\centering
\caption{Summary of main variables and hyperparameters used in the proposed method.}
\label{tab:notation_params}
\begingroup
\scriptsize
\setlength{\tabcolsep}{3pt}
\renewcommand{\arraystretch}{0.92}
\begin{tabular}{>{\raggedright\arraybackslash}p{0.2\textwidth}>{\raggedright\arraybackslash}p{0.30\textwidth}>{\raggedright\arraybackslash}p{0.46\textwidth}}
\toprule
Symbol & Meaning & Source or setting rule \\
\midrule
$\tau_d$ & Edge-length change threshold for dynamic topology removal & Fixed on the validation split; applied to all active THPG edges. \\
$\mathcal{W}_t(\cdot)$ & Graph-based non-rigid warp at frame $t$ & Interpolated from neighboring THPG-node displacements associated with each reference map point. \\
$\mathcal{A}_t$ & Index set of retained visible anchor nodes at frame $t$ & Obtained by filtering front-end THPG nodes using visibility, reprojection consistency, and tracking reliability. \\
$\Pi^{-1}(\cdot)$ & Camera back-projection operator & Converts a pixel or projected coordinate and its depth into a 3D point using current-view intrinsics. \\
$\tilde{k}(p)$ & Effective restoration weight in $E_{\mathrm{rest}}$ & Defined as $\tilde{k}(p)=k(p)r_{\mathrm{ssim}}(p)$, combining elastic-attribute propagation and appearance consistency. \\
$s_t,b_t$ & Frame-level scale and shift for pseudo-metric alignment & Updated during DDIM denoising; the denoising-step superscript is omitted for compact notation. \\
$\mathbf{p}_i$ & Current refined position of fused point $i$ & Initialized by $\mathbf{p}_i^0$ and updated during auxiliary post-fusion refinement. \\
$\lambda_e,\lambda_v$ & Front-end elastic and viscous regularization weights & Fixed scalar trade-off coefficients selected on the validation split. \\
$\lambda_g^0,\lambda_s^0,\lambda_r^0,\lambda_a^0,\lambda_b^0$ & Base weights of the back-end energy terms & Fixed on the validation split; $\lambda_g$, $\lambda_r$, and $\lambda_a$ are modulated by $c_t$ according to Eq.~(\ref{eq:scheduler}), while $\lambda_s$ and $\lambda_b$ remain fixed according to Eq.~(\ref{eq:scheduler-fixed}). \\
$\sigma,\gamma$ & 3D and 2D bandwidths in the anchor-edge weight $w_{ij}$ & Computed from the median 3D edge length and median projected 2D edge distance of active anchor edges in the current frame. \\
$N_{\mathrm{ref}}$ & Reference inlier count for tracking confidence & Set to 200 following NR-SLAM~\cite{bib9}; used only for confidence-normalization saturation, not as a hard failure threshold. \\
$\alpha_l$ & Latent guidance step size at DDIM step $l$ & Uses the DDIM test-time optimization learning-rate schedule; the base latent learning rate is reported in the experimental setup. \\
$\sigma_p$ & Distance bandwidth for refinement-edge weighting & Computed from the median reference edge length in $\mathcal{E}_{\mathrm{ref}}$. \\
$\lambda_{\mathrm{anc}},\lambda_{\mathrm{top}}$ & Anchor- and topology-preservation weights in post-fusion refinement & Fixed scalar trade-off coefficients selected on the validation split. \\
$\chi_i$ & Local anchor or fusion reliability of fused point $i$ & Inherited from anchor tracking consistency when available; otherwise estimated from local fusion reliability and normalized to $[0,1]$. \\
$\beta_i$ & Score-displacement strength in auxiliary refinement & Defined as $\beta_i=1-\chi_i$, so reliable anchor points are updated conservatively. \\
\bottomrule
\end{tabular}
\endgroup
\end{table*}
\setcounter{table}{9}

\clearpage
\begin{figure*}[p]
    \centering
    \includegraphics[width=\textwidth,height=0.82\textheight,keepaspectratio]{fig5.pdf}
    \caption{Qualitative comparison of dense depth completion under extremely sparse depth conditions. From top to bottom, the RGB image, sparse depth input, and depth predictions of different methods are shown. In the selected examples, several supervised depth completion methods show blurred boundaries, over-smoothing, confused depth ordering, and local collapse in medical images, whereas the proposed method shows improved boundary structure, local surface continuity, and overall scale consistency.}
    \label{fig:qual-backend}
\end{figure*}
\clearpage

\end{document}
